\chapter{Research Topic}

\section{Introduction}

The evolution of money is fundamentally a story about removing friction from human exchange. From the early days of bartering, where trade was heavily dependent on the double coincidence of wants, to the modern era of digital transactions, societies have continuously sought ways to make buying and selling faster, safer, and more convenient. In India, this evolution has recently taken a dramatic and unprecedented leap. For decades, the pinnacle of financial convenience for the middle and upper classes was the physical credit card---a rectangular piece of plastic that allowed consumers to "buy now, pay later." However, the introduction and subsequent explosion of the Unified Payments Interface (UPI) completely changed the rules of the game.

Initially, UPI was introduced merely as a fast, interoperable way to transfer money directly from a bank account using a smartphone and a simple virtual ID. It was revolutionary in its simplicity, but it lacked the crucial credit facility that many consumers relied upon for managing their monthly cash flows or financing large purchases. That changed drastically in 2022 when the Reserve Bank of India (RBI) allowed RuPay credit cards to be linked directly to UPI applications like Google Pay, PhonePe, and Paytm.

Suddenly, the traditional physical credit card faced its most significant existential threat. Consumers could now scan a simple QR code at a local vegetable vendor, an auto-rickshaw driver, or a neighborhood Kirana store and pay using their credit card limit, without ever taking the actual plastic card out of their wallet. This dissertation dives deep into this transition, exploring whether "Credit on UPI" is truly spelling the end for traditional physical credit cards in the Indian market, or if the two systems are destined to co-exist in a fragmented financial ecosystem.

\section{Historical Context and the Evolution of Payments}

To fully grasp the magnitude of the current shift, it is essential to trace the journey of payment systems. The concept of credit itself is not new; informal credit systems, such as the "Khata" or ledger system maintained by local shopkeepers in India, have existed for centuries. However, the formalization of consumer credit into a portable, universally accepted medium began in the mid-20th century in the United States with the introduction of the Diners Club card in 1950.

Initially made of cardboard, these cards eventually transitioned to celluloid and then to the ubiquitous Polyvinyl Chloride (PVC) plastic in the 1960s. Over the next six decades, global networks like Visa and Mastercard perfected the infrastructure required to authorize, clear, and settle transactions globally. The plastic card became synonymous with financial empowerment, status, and convenience.

Yet, this Western model of the "Four-Party" credit system (involving the consumer, merchant, acquiring bank, and issuing bank) was inherently expensive to maintain. It relied on costly physical infrastructure, primarily Point of Sale (POS) swipe machines, dedicated telephone or internet lines, and high transaction fees. Consequently, while credit cards flourished in developed nations, their penetration in price-sensitive, developing markets like India remained strictly limited to tier-1 cities, high-income individuals, and large organized retail outlets.

\section{The Indian Digital Payment Ecosystem and the Genesis of UPI}

India's digital payment ecosystem remained relatively stagnant until a series of structural shocks and technological innovations collided in the mid-2010s. The launch of the Jan Dhan Yojana in 2014 rapidly expanded bank account ownership, bringing millions into the formal financial sector. This was followed by the demonetization of high-value currency notes in 2016, which acted as a forced catalyst, pushing both consumers and reluctant merchants to adopt digital alternatives like mobile wallets almost overnight.

However, the true paradigm shift occurred with the launch of the Unified Payments Interface (UPI) by the National Payments Corporation of India (NPCI). Unlike mobile wallets, which were closed-loop systems requiring users to load money into a specific app, UPI was built as an open-architecture, interoperable public utility. It allowed instant, 24/7 bank-to-bank transfers using a simple Virtual Payment Address (VPA) or mobile number.

More importantly, from the merchant's perspective, UPI decimated the cost of acceptance. A merchant no longer needed a \rupee 10,000 POS machine; a simple printed QR code, which cost practically nothing to generate and display, was sufficient. This democratized digital payments, allowing even the smallest street vendor to accept digital money.

\section{The Convergence: Credit Meets UPI}

While UPI solved the problem of moving money, it did not solve the problem of creating credit. The credit card industry in India, though small compared to the overall population, was highly profitable and growing steadily, driven by aggressive marketing, co-branded cards, and lucrative reward programs.

This created a parallel financial existence. A consumer would use UPI for a \rupee 50 tea or a \rupee 200 cab ride, but would pull out their Visa or Mastercard for a \rupee 15,000 smartphone purchase or a restaurant dinner to earn reward points.

The RBI's decision to allow the linking of RuPay credit cards to UPI was a deliberate move to bridge this gap. By leveraging the massive, already-deployed network of over 650 million QR codes, the RBI aimed to expand the reach of formal credit beyond the limited 11 million POS terminals available in the country. This regulatory change set the stage for a direct confrontation between the established plastic card networks and the ascending QR-based ecosystem.

\section{Statement of the Problem}

Despite the obvious convenience of scanning a QR code to access credit, the transition is not frictionless. The core problem this dissertation addresses is the multifaceted resistance to the adoption of Credit on UPI, and the resulting impact on traditional credit cards.

Specifically, the problem involves:
1. \textbf{Consumer Habituation:} Assessing whether the convenience of mobile-based credit is sufficient to overcome the entrenched habit of swiping plastic, especially considering the differing reward structures.
2. \textbf{Merchant Viability:} Evaluating the severe pushback from small and medium merchants against the Merchant Discount Rate (MDR) applied to Credit on UPI transactions, which threatens their thin profit margins.
3. \textbf{Systemic Cannibalization:} Determining if Credit on UPI is actually replacing traditional card usage, or merely capturing a new segment of micro-credit transactions that were previously conducted in cash.

\section{Importance and Rationale of the Study}

This research is highly critical at this juncture for several compelling reasons:

First, from a \textbf{consumer behavior perspective}, the shift from physical plastic to digital QR codes represents a generational transformation. Understanding how different demographics, especially the digital-native Gen Z cohort compared to older generations, are adapting to this change helps forecast the future landscape of personal finance.

Second, the study exposes the \textbf{ground-level economic realities of merchants}. While macroeconomic reports celebrate the volume of digital transactions, they often overlook the microeconomic pain points of local shopkeepers. The friction caused by MDR fees is a significant hurdle that policymakers must address to ensure sustainable growth.

Third, this study provides vital insights for \textbf{banks, financial institutions, and legacy card networks}. As physical cards become potentially less relevant for daily use, banks are forced to fundamentally restructure their business models, rethink their reward programs, and redesign how credit limits are underwritten and distributed in a mobile-first, instant-gratification economy. This represents a multi-billion rupee industry shift.

\section{Scope and Limitations of the Dissertation}

\subsection{Scope}
The scope of this dissertation covers a comparative analysis between traditional physical credit cards (primarily Visa and Mastercard networks) and the emerging RuPay-based Credit on UPI system in India.

Geographically, the primary data collection and analysis are intensely focused on the Delhi National Capital Region (NCR). This area was chosen due to its high smartphone penetration, diverse demographic mix, and vibrant retail sector encompassing everything from premium malls to unregulated street markets.

The study targets a diverse demographic of consumers across various age groups (Gen Z, Millennials, and Gen X) and a mix of merchants. It covers domestic retail transactions and consumer credit behavior over the period from 2020 to 2025, providing a comprehensive view backed by primary survey data and trend analysis from secondary sources like the RBI and NPCI.

\subsection{Limitations}
While extensive, this study acknowledges certain limitations:
1. \textbf{Geographical Constraints:} The primary data is limited to the Delhi NCR region. While representative of urban Indian behavior, it may not fully capture the nuances of rural or tier-3 city adoption where credit penetration is historically much lower.
2. \textbf{Evolving Regulatory Landscape:} The regulations surrounding Credit on UPI, particularly regarding MDR caps and network interoperability, are highly dynamic. Changes implemented by the RBI post-2025 may alter the long-term trajectory predicted in this study.
3. \textbf{Self-Reporting Bias:} As the primary data relies on self-reported surveys, there is a natural possibility of discrepancies between what respondents claim they do and their actual financial behavior.
